\subsection{Comparison with the opened semigroup literature}
We opened Sweeney's Theorem 2.7 and Corollary 2.8, and the
integral projective-resolution method in Sections 3.1--3.5,
including Examples 3.1 and 3.3\cite{CP2-Sweeney2025}.
The idempotent-ideal projectives in our certificates implement
the same homological principle: exact embedded kernel lattices
are covered by direct summands of free modules, trivial
coinvariants turn the maps into integer matrices, and Smith
normal form records torsion. Our independent verifier checks
the full ideals, actions, differential equations and saturated
images, rather than assuming exactness from a truncated rank
calculation. The degree-two bar comparisons and cup primitives
are additional data needed to distinguish a projective plane
from a space having only the same additive homology. We have
not run Sweeney's external software package, and do not claim
equality with an unreceived output catalog.

Fiedorowicz's unnumbered lemma on pages 33--34 of his 2002
paper gives the five-element rectangular-band sphere monoid
and a projective right-module resolution\cite{CP2-Fiedorowicz2002}.
Our rectangular-band calculation contains this example, with
the left-module orientation used in Sweeney's Example 3.1.
We opened that entire three-page paper, including its separate
unnumbered simplicial-monoid counterexample theorem on pages
34--35. That theorem is not a minimality statement for finite
projective-plane monoids.

We also opened the introduction and Theorems A--C of
Martinian--Steinberg\cite{CP2-MartinianSteinberg2025}.
Their finite wedges of simply connected Moore spaces do not
supply the projective-plane cup square: positive-degree
products on such a wedge vanish. The proof-level comparison
with that paper remains explicitly unfinished. Exact source
versions, URLs, hashes, read locators and remaining obligations
are recorded in \path{notes/root_source_record_R11.json} and
the report. Availability of an attached source is not treated
as having checked all its proofs.
